% =====================================================================
% Group A -- Orbital mechanics, chapters 2, 3, 6 (the "third" orbital question)
% Source: latex_notes/SFM_02_*, SFM_03_*, SFM_06_*
% =====================================================================
\qgroup{A}{Orbital mechanics -- Chapters 2, 3, 6}

At the exam, one of the three orbital-mechanics questions is on a chapter other than 5 and 8.
The questions below are ordered by topic: Keplerian trajectories (A1--A4), impulsive transfers (A5--A7),
rocket dynamics (A8--A11). Adjacent questions share material but each answer is self-contained.

\begin{center}
\small
\begin{tabular}{@{}lllr@{}}
\toprule
Q & Topic & Asked & Page\\
\midrule
\ref{q:A1} & First integral of angular momentum & Feb 2025 & \pageref{q:A1}\\
\ref{q:A2} & Classification of Keplerian orbits ($a$, $e$) & list & \pageref{q:A2}\\
\ref{q:A3} & Horizontal velocity at given altitude & list & \pageref{q:A3}\\
\ref{q:A4} & Geosynchronous and geostationary orbits, ground tracks & Jan 2023 & \pageref{q:A4}\\
\ref{q:A5} & Impulsive-thrust approximation & list & \pageref{q:A5}\\
\ref{q:A6} & Hohmann vs bielliptic transfer & $\times$2, Jan 2025, Jun 2025 & \pageref{q:A6}\\
\ref{q:A7} & Ellipse $\to$ hyperbola: two optimal strategies & list & \pageref{q:A7}\\
\ref{q:A8} & $\Delta V$-loss integrals & $\times$2, Sep 2024 & \pageref{q:A8}\\
\ref{q:A9} & Propulsive velocity and losses vs trajectory & Jul 2025 & \pageref{q:A9}\\
\ref{q:A10} & Staging: mass parameters, why staging & list & \pageref{q:A10}\\
\ref{q:A11} & Tsiolkovsky and multistage $\Delta V$ & Sep 2025 & \pageref{q:A11}\\
\bottomrule
\end{tabular}
\end{center}

% ---------------------------------------------------------------------
\question{A1}{First integral of angular momentum}
% source: cap02.tex l.168-201 (h), l.589-615 (theta-dot), l.797-848 (circular), cap03 l.51-75 (impulses)
\begin{qbox}
Prove the first integral of angular momentum for a Keplerian orbit and give some application examples.
\qmeta{Feb 2025 (listed as a 2025 question).}
\end{qbox}

\begin{outline}
\begin{enumerate}
\item Restricted two-body problem: equation of motion $\ddot{\underline r}=-\frac{\mu}{r^2}\hat r$ (central force).
\item Definition $\underline h:=\underline r\times\underline v$ and proof $\dot{\underline h}=\underline 0$.
\item Consequence (A): planar motion, $\hat h$ fixed $\Rightarrow$ orbit plane fixed in inertial space.
\item Consequence (B): Kepler's 2nd law $h=2\,dA/dt$ (+ sketch).
\item Applications: $\dot\theta_*=h/r^2$, $v_\theta=h/r$, $r_pv_p=r_Av_A$, $p=h^2/\mu$, orbit plane ($i,\Omega$), $h=0$ rectilinear,
impulses ($h^+=r\,v_\theta^+$), perturbations ($\dot{\underline h}=\underline r\times\underline f$).
\end{enumerate}
\end{outline}

\answer
\paragraph{Setting.} In the restricted two-body problem ($m_1\gg m_2$, so $m_1+m_2\simeq m_1$ and the effect of the
spacecraft on the main body is neglected) the position $\underline r$ of the spacecraft relative to the centre of the
attracting body obeys
\begin{equation}
\frac{d^{2}\underline{r}}{dt^{2}}=-\frac{\mu}{r^{3}}\underline{r}=-\frac{\mu}{r^{2}}\hat{r},
\qquad \mu=Gm_1 \ (\mu_\oplus=398600.4\ \mathrm{km^3/s^2}).
\end{equation}
The force is \emph{central}: it is always aligned with $\underline r$. A \emph{first integral} is a quantity that remains
constant in time along the motion.

\paragraph{Definition and proof.} The specific angular momentum (per unit mass) is
\begin{equation}
\underline{h}:=\underline{r}\times\underline{v},\qquad \underline{v}=\frac{d\underline{r}}{dt}.
\end{equation}
Differentiating in the inertial frame and using the equation of motion:
\begin{equation}
\frac{d\underline{h}}{dt}=\underbrace{\frac{d\underline{r}}{dt}\times\underline{v}}_{\underline v\times\underline v=\underline 0}
+\underline{r}\times\underbrace{\frac{d^{2}\underline{r}}{dt^{2}}}_{-\frac{\mu}{r^{2}}\hat{r}}
=\underline 0-\frac{\mu}{r^{2}}\,\underbrace{\underline r\times\hat r}_{\underline 0}=\underline 0
\;\Rightarrow\;\underline{h}=\text{constant}.
\end{equation}
The first term vanishes because a vector crossed with itself is zero; the second because the gravitational acceleration
is parallel to $\underline r$. Hence $\underline h$ is a vector first integral (3 scalar constants), fixed by the initial
conditions: $\underline h=\underline r_0\times\underline v_0$.

\paragraph{Consequence (A): the trajectory is planar.} Since $\underline h=\text{const}$, also $\hat h=\text{const}$
(if $h\neq 0$). By definition $\underline r\cdot\underline h=0$ and $\underline v\cdot\underline h=0$ at all times, so
$\underline r$ and $\underline v$ always lie in the plane through the attracting body orthogonal to $\hat h$: the orbit
plane is fixed in inertial space. (This is why the orbit plane is described by two constant angles, $i$ and $\Omega$.)

\paragraph{Consequence (B): equal areas in equal times (Kepler's 2nd law).} In $dt$ the position vector sweeps the
triangle of sides $\underline r$ and $d\underline r=\underline v\,dt$, of area $dA=\frac12|\underline r\times d\underline r|$:
\begin{equation}
h\,dt=|\underline{r}\times\underline{v}|\,dt=|\underline{r}\times d\underline{r}|=2\,dA
\;\Longrightarrow\; h=2\frac{dA}{dt}=\text{constant},
\end{equation}
i.e.\ the areolar velocity $dA/dt$ is constant.

\begin{center}
\includegraphics[width=0.42\linewidth]{kepler2}\\
{\footnotesize Equal times $\Rightarrow$ equal areas: the spacecraft is fastest near periapsis.}
\end{center}

\paragraph{Application examples.}
\begin{enumerate}[label=(\arabic*)]
\item \textbf{Angular rate along the orbit.} Writing $\underline v=\dot r\,\hat r+\underline\omega\times\underline r$ with
$\underline\omega=\dot\theta_*\hat h$: $\underline h=\underline r\times(\underline\omega\times\underline r)=\underline\omega r^2$, so
\begin{equation}
\dot\theta_*=\frac{h}{r^2}=\sqrt{\frac{\mu}{p^3}}\left(1+e\,c_{\theta_*}\right)^2 ,
\end{equation}
the differential equation that relates true anomaly and time (basis of Kepler's equation): the motion is fastest at periapsis.
\item \textbf{Horizontal velocity.} $h=r\,v_\theta$ $\Rightarrow$ $v_\theta=h/r$. At the apses ($v_r=0$):
$r_p v_p=r_A v_A=h$, hence $v_p>v_A$ (consistent with the 2nd law).
\item \textbf{Shape of the orbit.} Combined with the eccentricity vector $\underline e=-\hat r+\underline v\times\underline h/\mu$
it gives the polar equation $r=\dfrac{h^2/\mu}{1+e\,c_{\theta_*}}$: the semilatus rectum is $p=h^2/\mu$.
\item \textbf{Orientation of the orbit plane.} $\hat h$ gives the inclination $c_i=\hat h\cdot\hat c_3$ and, with the node
line $\hat c_3\times\underline h$, the RAAN $\Omega$.
\item \textbf{Circular orbits.} $r=R$ and $h=$ const $\Rightarrow$ $\underline\omega=\underline h/R^2$ constant: uniform motion.
\item \textbf{Rectilinear trajectories.} $h=0$ iff $\underline v\parallel\underline r$: $p=0$, $e=1$, the motion is along a line
through the attracting body.
\item \textbf{Impulsive manoeuvres (Ch.\,3).} After an impulse at radius $r$: $h^+=r\,v_\theta^+$, then
$e^+=\sqrt{1-h^{+2}/(\mu a^+)}$; in the Hohmann proof $v_+c_\phi=h^+/R_k=\sqrt{\mu p^+}/R_k$.
\item \textbf{Perturbed motion (Ch.\,8).} With a perturbing acceleration $\underline f$ the integral is lost:
$\dot{\underline h}=\underline r\times\underline f$. Only the out-of-plane component $f_h$ rotates the plane (changes $i$,
$\Omega$); e.g.\ third-body perturbations make $\underline h$ precess.
\end{enumerate}

\begin{keyf}
$\underline h=\underline r\times\underline v$, \quad $\dfrac{d\underline h}{dt}=\underline v\times\underline v+\underline r\times\left(-\dfrac{\mu}{r^2}\hat r\right)=\underline 0$,
\quad $h=2\,\dfrac{dA}{dt}$, \quad $\dot\theta_*=\dfrac{h}{r^2}$, \quad $p=\dfrac{h^2}{\mu}$, \quad $h=r\,v_\theta=r_pv_p=r_Av_A$.
\end{keyf}

\pitfalls
\begin{itemize}
\item State the hypothesis explicitly: \emph{central} force (restricted two-body problem, spherical attracting body).
The proof fails for any non-central perturbation.
\item Say \emph{why} each term vanishes ($\underline v\times\underline v=\underline 0$; $\underline r\parallel\ddot{\underline r}$).
\item $h$ is constant as a \emph{vector}: magnitude (2nd law) and direction (planar motion). Mention both consequences.
\item $\underline h$ and $\underline e$ give 5 independent scalar integrals, not 6, because $\underline h\cdot\underline e=0$.
\end{itemize}

\further
The angular momentum integral is the orbital analogue of the conservation of angular momentum of a system subject to
central (internal) forces: in the $N$-body problem the total $\underline H_c$ about the centre of mass is conserved and
defines the invariant Laplace plane (\qref{B6}).

\source{cap02.tex, \S First Integrals (l.168--201), \S Position in time (l.589--615), \S Circular orbits (l.797--848);
cap03.tex l.51--75; cap08.tex l.87--116. Formulario: \fml{2-HEV}, \fml{2-TDT}.}

% ---------------------------------------------------------------------
\question{A2}{Classification of Keplerian orbits in terms of $a$ and $e$}
% source: cap02.tex l.253-300 (cartesian eq.), 300-371 (conics), 554-587 (energy), 881-927 (rectilinear, classification)
\begin{qbox}
Describe the general classification of Keplerian orbits in terms of semi-major axis and eccentricity.
\qmeta{list by chapter (Ch.\,2).}
\end{qbox}

\begin{outline}
\begin{enumerate}
\item First integrals $\underline h$, $\underline e$ $\Rightarrow$ polar equation $r=p/(1+e\,c_{\theta_*})$, $p=h^2/\mu$.
\item Cartesian form $x^2(1-e^2)+y^2+2pex=p^2$: conic section, the attracting body is at a focus.
\item Definition of $a=p/(1-e^2)$ for all conics ($a>0$, $a\to\infty$, $a<0$).
\item Energy $\mathcal E=-\mu/(2a)$: sign of $\mathcal E$ $\leftrightarrow$ type of orbit.
\item Degenerate case $h=0$: rectilinear trajectories ($e=1$), three types.
\item Summary table $\mathcal E$ (or $a$) vs $e$; sketch of the conics.
\end{enumerate}
\end{outline}

\answer
\paragraph{From the first integrals to the conics.} In the restricted two-body problem $\underline h=\underline r\times\underline v$
and $\underline e=-\hat r+\underline v\times\underline h/\mu$ are constant. Projecting $\underline e$ on $\underline r$
($\theta_*$ = true anomaly, angle from $\underline e$ to $\underline r$):
\begin{equation}
r\,e\cos\theta_*=-r+\frac{h^2}{\mu}\;\Longrightarrow\; r=\frac{p}{1+e\cos\theta_*},\qquad p:=\frac{h^2}{\mu}\ \text{(semilatus rectum)}.
\end{equation}
In a frame centred in the attracting body with $x\parallel\hat e$, using $x=r\,c_{\theta_*}$, $y=r\,s_{\theta_*}$:
\begin{equation}
x^{2}\left(1-e^{2}\right)+y^{2}+2p\,e\,x=p^{2}
\end{equation}
which is a conic section with the attracting body at one focus (Kepler's 1st law):
\begin{itemize}
\item $0\le e<1$: \textbf{ellipse} (circle with the body at the centre if $e=0$);
\item $e=1$: \textbf{parabola} ($y^2+2px=p^2$, $r_p=p/2$);
\item $e>1$: \textbf{hyperbola} (one branch, the one around the focus).
\end{itemize}

\paragraph{Semi-major axis.} Completing the square gives, for the ellipse, semi-axes $a=\dfrac{p}{1-e^2}$,
$b=\dfrac{p}{\sqrt{1-e^2}}$, focal distance $c=ae$. By \emph{convention} the same definition is kept for hyperbolas,
\begin{equation}
a=\frac{p}{1-e^{2}}\quad\Rightarrow\quad a>0\ (e<1),\qquad a\to\infty\ (e=1),\qquad a<0\ (e>1),
\end{equation}
so that $r_p=a(1-e)$ holds for every conic and $r_A=a(1+e)$ for the ellipse.

\paragraph{Energy.} The specific energy $\mathcal E=\frac{v^2}{2}-\frac{\mu}{r}$ is a further integral. Inserting $r(\theta_*)$ and
the speed $v=\sqrt{\mu/p}\,\sqrt{1+e^2+2e\,c_{\theta_*}}$, the dependence on $\theta_*$ cancels:
\begin{equation}
\mathcal E=-\frac{\mu}{2p}\left(1-e^{2}\right)=-\frac{\mu}{2a}.
\end{equation}
Hence the sign of $\mathcal E$ (equivalently of $a$) classifies the orbit:
\begin{center}
\small
\begin{tabular}{@{}lllll@{}}
\toprule
orbit & $e$ & $a$ & $\mathcal E$ & potential vs kinetic energy\\
\midrule
circle & $0$ & $=r$ & $<0$ & $|\mathrm{PE}|=2\,\mathrm{KE}$\\
ellipse (closed, periodic) & $0<e<1$ & $>0$ & $<0$ & $|\mathrm{PE}|>\mathrm{KE}$\\
parabola (escape, $v_\infty=0$) & $1$ & $\infty$ & $=0$ & $|\mathrm{PE}|=\mathrm{KE}$\\
hyperbola (escape, $v_\infty>0$) & $>1$ & $<0$ & $>0$ & $|\mathrm{PE}|<\mathrm{KE}$\\
\bottomrule
\end{tabular}
\end{center}
Characteristic quantities: ellipse $r_{p,A}=a(1\mp e)$, $T=2\pi\sqrt{a^3/\mu}$; parabola $v=\sqrt{2\mu/r}\to0$ at infinity;
hyperbola $v_\infty=\sqrt{-\mu/a}$, true anomaly limited to $|\theta_*|<\theta_*^M=\arccos(-1/e)$, deflection
$\xi=2\theta_*^M-\pi$.

\paragraph{Degenerate case: rectilinear trajectories ($h=0$).} If $\underline v\parallel\underline r$ then $\underline h=\underline 0$,
$e=|-\hat r|=1$, $p=0$, $r_p=0$: the motion is along a straight line through the attracting body ($\theta_*=\pm\pi$). Three
types, classified again by $a$ (or $\mathcal E$):
$a>0,\ \mathcal E<0$ \emph{rectilinear ellipse} (segment: the body falls back);
$a\to\infty,\ \mathcal E=0$ \emph{rectilinear parabola} ($v\to0$ as $r\to\infty$);
$a<0,\ \mathcal E>0$ \emph{rectilinear hyperbola} ($v\not\to0$).

\paragraph{Summary.} A Keplerian path is identified by \textbf{$a$ (or $\mathcal E$) and $e$}:
\begin{center}
\small
\begin{tabular}{@{}c|c|c|c@{}}
\toprule
$\mathcal E$ & $0\le e<1$ & $e=1$ & $e>1$\\
\midrule
$>0$ ($a<0$) & -- & rectilinear hyperbola & hyperbola\\
$=0$ ($a\to\infty$) & -- & parabola, rectilinear parabola & --\\
$<0$ ($a>0$) & ellipse & rectilinear ellipse & --\\
\bottomrule
\end{tabular}
\end{center}
Conic sections have $h\neq0$, $e\ge0$; rectilinear trajectories have $h=0$, $e=1$. Given $r$ and $v$ at any point:
$a=\left(\frac{2}{r}-\frac{v^2}{\mu}\right)^{-1}$ (vis-viva) and $e=\sqrt{1-\frac{h^2}{\mu a}}$.

\begin{center}
\includegraphics[width=0.62\linewidth]{conics}\\
{\footnotesize Conics with the same periapsis radius $r_p$: $e$ increases from the circle to the hyperbola.}
\end{center}

\begin{fig2draw}
Focus $F$ and a common periapsis; circle, an ellipse closing behind, the parabola (arms becoming parallel) and the
hyperbola (arms approaching the asymptotes, more open than the parabola). Write $e$ and the sign of $a$, $\mathcal E$ on each.
\end{fig2draw}

\begin{keyf}
$r=\dfrac{p}{1+e\,c_{\theta_*}}$, \quad $p=\dfrac{h^2}{\mu}=a(1-e^2)$, \quad $\mathcal E=\dfrac{v^2}{2}-\dfrac{\mu}{r}=-\dfrac{\mu}{2a}$,
\quad $e=\sqrt{1-\dfrac{h^2}{\mu a}}$, \quad $v_\infty=\sqrt{-\mu/a}$.
\end{keyf}

\pitfalls
\begin{itemize}
\item The convention $a<0$ for hyperbolas: forgetting it gives $r_p<0$ or imaginary $v_\infty$.
\item Do not forget the \emph{rectilinear} (degenerate) trajectories: they share $e=1$ with the parabola but have any $\mathcal E$.
\item The circle is the special ellipse $e=0$ (attracting body at the centre), not a separate class.
\item $e$ alone does not identify the orbit size; $a$ alone does not identify the shape: both are needed.
\end{itemize}

\source{cap02.tex l.230--300 (polar and Cartesian equations), l.300--371 (ellipse, parabola, hyperbola), l.554--587 (energy),
l.881--927 (rectilinear trajectories, classification table), l.929--945 (vis-viva). Formulario: \fml{2-POL}, \fml{2-ENE}, \fml{2-VIS}.}

% ---------------------------------------------------------------------
\question{A3}{Trajectories from a horizontal velocity at a given altitude}
% source: cap02.tex l.532-552 (cosmic velocities), l.929-1005 (cases a-g), l.870-879 (ballistic)
\begin{qbox}
A spacecraft is at a specified altitude above Earth and has horizontal velocity. Describe the different Keplerian
trajectories it can follow for various velocity magnitudes.
\qmeta{list by chapter (Ch.\,2).}
\end{qbox}

\begin{outline}
\begin{enumerate}
\item Data: $r=R_E+$altitude, $\underline v_0\perp\underline r$ $\Rightarrow$ $P_0$ is an apse ($v_r=0$).
\item Cosmic velocities $v_I=\sqrt{\mu/r}$ (circular), $v_{II}=\sqrt{2\mu/r}$ (escape).
\item $a$ from vis-viva, $e$ from $h=rv_0$: $e=|v_0^2/v_I^2-1|$.
\item Seven cases (a)--(g) with sketch; $P_0$ apoapsis below $v_I$, periapsis above.
\item Remark: low $v_0$ $\Rightarrow$ periapsis inside the Earth: ballistic arc (impact).
\end{enumerate}
\end{outline}

\answer
\paragraph{Data and cosmic velocities.} Let $P_0$ be at distance $r$ from the Earth's centre ($r=R_E+\text{altitude}$) and let
the velocity $\underline v_0$ be horizontal, i.e.\ orthogonal to $\underline r$. Two \emph{cosmic velocities} are defined at $P_0$:
\begin{equation}
v_I=\sqrt{\frac{\mu}{r}}\ \ \text{(1st: circular orbit of radius $r$)},\qquad
v_{II}=\sqrt{\frac{2\mu}{r}}=\sqrt2\,v_I\ \ \text{(2nd: escape, parabola)}.
\end{equation}

\paragraph{Orbit elements as functions of $v_0$.} With $\underline v_0\perp\underline r$: $v_r=0$, so $P_0$ is an apse (periapsis or
apoapsis), and $h=r\,v_0$. Vis-viva and angular momentum give
\begin{equation}
a=\frac{1}{\dfrac{2}{r}-\dfrac{v_0^{2}}{\mu}},\qquad p=\frac{h^2}{\mu}=r\,\frac{v_0^2}{v_I^2},\qquad
e=\left|\frac{p}{r}-1\right|=\left|\frac{v_0^{2}}{v_I^{2}}-1\right|,
\end{equation}
the last one because at an apse $r=p/(1\pm e)$. With $k=v_0^2/v_I^2$ the \emph{other} apse is at $r'=r\,k/(2-k)$.
As $v_0$ grows, $\mathcal E=\frac{v_0^2}{2}-\frac{\mu}{r}$ grows monotonically.

\paragraph{The cases.}
\begin{enumerate}[label=(\alph*)]
\item $v_0=0$: \textbf{rectilinear ellipse}, straight fall towards the centre ($h=0$, $e=1$, $a=r/2$).
\item $0<v_0<v_I$: \textbf{ellipse with $P_0$ as apoapsis} ($e=1-k$); the periapsis radius $r_p=r\,k/(2-k)$ increases with $v_0$.
If $r_p<R_E$ the ellipse intersects the Earth's surface: the path is a \emph{ballistic trajectory} (arc of ellipse with
periapsis inside the Earth, as for ICBMs) and the spacecraft re-enters/impacts.
\item $v_0=v_I$: \textbf{circular orbit} of radius $r$ ($e=0$, $a=r$).
\item $v_I<v_0<v_{II}$: \textbf{ellipse with $P_0$ as periapsis} ($e=k-1$); the apoapsis radius $r_A=r\,k/(2-k)$ increases
with $v_0$ and tends to infinity as $v_0\to v_{II}$.
\item $v_0=v_{II}$: \textbf{parabola} with periapsis $r_p=r$ ($e=1$, $\mathcal E=0$): escape with zero velocity at infinity.
\item $v_0>v_{II}$: \textbf{hyperbola} with periapsis at $P_0$; $e$ increases with $v_0$, $v_\infty=\sqrt{v_0^2-v_{II}^2}$.
\item $v_0\to\infty$: $e\to\infty$, the path tends to the \textbf{straight line} orthogonal to $\underline r$ (gravity has no time to bend it).
\end{enumerate}

\begin{center}
\includegraphics[width=0.86\linewidth]{vhoriz}\\
{\footnotesize Paths computed with $r=p/(1+e\cos\theta_*)$ for $v_0/v_I=0.7,\ 1,\ 1.25,\ \sqrt2,\ 1.75$.}
\end{center}

\begin{fig2draw}
Earth centre $B$, point $P_0$ to its right, short arrow $\underline v_0$ upwards. Draw: small ellipse inside (apoapsis at $P_0$),
circle through $P_0$, larger ellipse outside (periapsis at $P_0$), parabola and hyperbola opening to the left, and
the radial segment $BP_0$ for $v_0=0$. Label each with the velocity interval.
\end{fig2draw}

\begin{keyf}
$v_I=\sqrt{\mu/r}$, \ $v_{II}=\sqrt{2\mu/r}$, \ $a=\left(\dfrac2r-\dfrac{v_0^2}{\mu}\right)^{-1}$, \ $e=\left|\dfrac{v_0^2}{v_I^2}-1\right|$,
\ $h=rv_0$, \ other apse $r'=r\dfrac{k}{2-k}$, $k=\dfrac{v_0^2}{v_I^2}$.
\end{keyf}

\pitfalls
\begin{itemize}
\item Below $v_I$ the starting point is the \emph{apoapsis}, above $v_I$ it is the \emph{periapsis}.
\item $v_{II}=\sqrt2\,v_I$ holds at the same $r$; the escape speed is not "infinite energy".
\item Mention the physical limit: for small $v_0$ the periapsis falls inside the Earth (ballistic arc).
\end{itemize}

\source{cap02.tex l.532--552 (cosmic velocities), l.870--879 (ballistic trajectories), l.929--1005 (cases (a)--(g)).
The closed-form $e=|v_0^2/v_I^2-1|$ and $r'$ follow from the notes' formulas. Formulario: \fml{2-COS}, \fml{2-V0C}.}

% ---------------------------------------------------------------------
\question{A4}{Geosynchronous and geostationary orbits; ground tracks}
% source: cap02.tex l.850-868 (geosynchronous), l.1657-1843 (ground track), cap08 l.80-85 (J22 resonance)
\begin{qbox}
Define geosynchronous and geostationary orbits, commenting on their properties. Provide a graphical representation of the
ground track for zero/non-zero inclination and zero/non-zero eccentricity.
\qmeta{Jan 2023.}
\end{qbox}

\begin{outline}
\begin{enumerate}
\item Sidereal day $T_{sid}=2\pi/\omega_E=86164$\,s; geosynchronous: $T=T_{sid}$ $\Rightarrow$ $a=42164$\,km, any $e$, $i$.
\item Geostationary = geosynchronous + $e=0$ + $i=0$: fixed sub-satellite point on the equator.
\item Ground track definition: $\phi(t)$, $\lambda_g(t)=\xi-\theta_g$; $\phi=\arcsin(s_{\theta_t}s_i)$, latitude band $\pm i$.
\item Four sketches: ($i=0,e=0$) point; ($i\neq0,e=0$) figure eight; ($i=0,e\neq0$) segment on the equator;
($i\neq0,e\neq0$) distorted eight/drop, depending on $\omega$.
\item Properties: closed track for direct orbits (repeats every day); symmetries; motion E/W; uses and perturbations.
\end{enumerate}
\end{outline}

\answer
\paragraph{Definitions.} The \emph{sidereal day} is the time the Earth takes for one rotation with respect to an inertial frame:
$T_{sid}=2\pi/\omega_E=86164$\,s ($\omega_E=7.292115\cdot10^{-5}$\,rad/s), shorter than the solar day (86400\,s).
An orbit is \textbf{geosynchronous} if its period equals one sidereal day:
\begin{equation}
2\pi\sqrt{\frac{a^{3}}{\mu_{E}}}=T_{\text{sid}}\;\Longrightarrow\;
a_{gs}=\left[\mu_{E}\left(\frac{T_{\text{sid}}}{2\pi}\right)^{2}\right]^{1/3}=42164\ \text{km}
\end{equation}
(altitude $\approx35786$\,km for a circular orbit). \emph{No assumption is made on $i$ and $e$}, apart from a periapsis high
enough to avoid the atmosphere.

An orbit is \textbf{geostationary} if it is at the same time (1) geosynchronous ($a=42164$\,km), (2) circular ($e=0$),
(3) equatorial ($i=0$). The satellite then moves with the same angular velocity as the Earth, in the same plane and
direction: the sub-satellite point is a \textbf{fixed point on the equator} and the satellite appears motionless to a
ground observer. Geostationary orbits are a subset of geosynchronous orbits.

\paragraph{Ground track.} It is the locus of the sub-satellite points (projection of the satellite on the rotating Earth),
plotted on a Mercator map as latitude $\phi(t)$ vs geographical longitude $\lambda_g(t)$:
\begin{equation}
s_\phi=s_{\theta_t}s_i,\qquad \lambda_g=\xi-\theta_g,\qquad \theta_g=\theta_{g0}+\omega_E(t-t_0),
\end{equation}
with $\theta_t=\omega+\theta_*$ (argument of latitude), $\xi$ the absolute (inertial) longitude and $\theta_g$ the
longitude of Greenwich. For direct orbits $-i\le\phi\le i$.

\paragraph{The four cases} (tracks computed with Kepler's equation, $a=42164$\,km):
\begin{itemize}
\item \textbf{$i=0$, $e=0$} (geostationary): a single \textbf{point} on the equator.
\item \textbf{$i\neq0$, $e=0$}: a symmetric \textbf{figure eight} between latitudes $\pm i$, centred on the equator, with the
crossing point at the equator. The satellite moves at constant inertial speed, but its eastward component varies with
latitude, so it alternately leads and lags the Earth (small longitude excursion, $\propto i^2$). Symmetric w.r.t.\ the
meridian through $\phi_{\max}$ and w.r.t.\ the equator crossing (both symmetries, $e=0$).
\item \textbf{$i=0$, $e\neq0$}: a \textbf{segment on the equator} travelled back and forth once a day: near perigee the satellite
is faster than the Earth (moves east), near apogee slower (moves west); the amplitude grows with $e$ (about $\pm2e$\,rad
for small $e$ \notinnotes).
\item \textbf{$i\neq0$, $e\neq0$}: a \textbf{distorted (asymmetric) eight} or a \textbf{drop/loop} shape, still closed. The shape depends on
$\omega$: with $\omega=0,\pi$ the track is symmetric w.r.t.\ the points where it crosses the equator; with $\omega=\pm\pi/2$
it is symmetric w.r.t.\ the meridian through the maximum and minimum latitudes (apocentre/pericentre).
\end{itemize}
\begin{center}
\includegraphics[width=0.95\linewidth]{geosync}
\end{center}

\paragraph{Properties.}
\begin{itemize}
\item Since the period is one sidereal day, the motion is synchronous with the Earth's rotation and the \textbf{track repeats
every day}. For direct orbits ($0<i<\pi/2$) the track is \textbf{closed} and spans a limited range of longitudes; it is travelled
partly towards East and partly towards West. For retrograde geosynchronous orbits ($\pi/2<i<\pi$) the track is repeating
but covers all longitudes (always travelled westward).
\item General rule: after one period a direct track drifts east if $T<T_{sid}$, west if $T>T_{sid}$; geosynchronous is the
boundary case with zero drift.
\item Geostationary orbits give continuous coverage of about one third of the Earth with fixed antennas
(telecommunications, meteorology); they cannot see the polar regions.
\item Perturbations: at GEO the dominant one is the sectoral harmonic $J_{22}$ (ellipticity of the equator), amplified by
resonance because the satellite is fixed with respect to the Earth; Sun and Moon make $i$ grow, SRP changes $e$
$\Rightarrow$ station keeping is needed.
\end{itemize}

\begin{fig2draw}
Four small $\lambda_g$--$\phi$ frames: a dot at the origin; an "8" between $\pm i$; a horizontal segment on the equator with a
double arrow; a lopsided "8" (or a drop) for $e\neq0$, with a note on the symmetry given by $\omega$.
\end{fig2draw}

\begin{keyf}
$T_{sid}=\dfrac{2\pi}{\omega_E}=86164$\,s, \quad $a_{gs}=\left[\mu_E\left(\dfrac{T_{sid}}{2\pi}\right)^2\right]^{1/3}=42164$\,km,
\quad $s_\phi=s_{\theta_t}s_i$, \quad $\lambda_g=\xi-\theta_g$.
\end{keyf}

\pitfalls
\begin{itemize}
\item Use the \emph{sidereal} day (86164\,s), not 24\,h.
\item "Geosynchronous" does not imply circular or equatorial; "geostationary" requires all three conditions.
\item The figure eight of the circular inclined case is symmetric and centred on the equator; the eccentric cases are not.
\end{itemize}

\source{cap02.tex l.850--868 (geosynchronous orbits), l.1657--1739 (ground track), l.1741--1796 (latitude limits, symmetry),
l.1798--1855 (geosynchronous and geostationary tracks, motion along the track), l.1857--1960 (examples 1--3);
cap08.tex l.80--85 ($J_{22}$ at GEO). Formulario: \fml{2-GEO}, \fml{2-GTK}, \fml{2-GTD}.}

% ---------------------------------------------------------------------
\question{A5}{Impulsive-thrust approximation}
% source: cap03.tex l.18-113
\begin{qbox}
Describe the impulsive-thrust approximation -- its validity assumptions and the change in position and velocity due to an
impulsive $\Delta V$. Give one mission scenario where the approximation is valid and one where it is not.
\qmeta{list by chapter (Ch.\,3).}
\end{qbox}

\begin{outline}
\begin{enumerate}
\item Idea: high thrust for a short time $\Rightarrow$ instantaneous velocity change; limit $T\to\infty$, $\Delta t\to0$, $T\Delta t=\Delta v$ finite.
\item Effects: $\underline r^+=\underline r^-$, $\underline v^+=\underline v^-+\Delta\underline v$ (magnitude and direction).
\item Validity: burn time $\ll$ orbital period, negligible displacement and gravity/drag effects during the burn; optimistic estimate.
\item New orbit: $a^+$ from energy, $e^+$ from $h$, $\theta_*^+$ from $r$ and $v_r$; apsidal line rotates.
\item Examples: valid (chemical orbit transfers, Hohmann, injection); not valid (launcher ascent, low-thrust).
\end{enumerate}
\end{outline}

\answer
\paragraph{The approximation.} If a \textbf{high thrust} acts for a \textbf{short time} $\Delta t$, its effect can be approximated as an
\emph{instantaneous} change of velocity, while the position does not change. Mathematically (thrust per unit mass $T$)
\begin{equation}
T\,\Delta t=\Delta v,
\end{equation}
and the impulsive approximation is the limit $T\to\infty$, $\Delta t\to0$ with the product $T\Delta t=\Delta v$ finite. Through a
\textbf{velocity impulse}: (a) the velocity changes in direction and magnitude; (b) the position does not change.

\paragraph{Validity assumptions.}
\begin{itemize}
\item High-thrust propulsion (chemical engines): the burn duration is very short compared with the orbital period, so the
spacecraft moves negligibly during the burn and the change of $\underline r$ is negligible.
\item During the burn gravity (and drag) do not change the velocity appreciably: no gravity/drag losses; Tsiolkovsky's law
can be used to compute the propellant.
\item The result is an \emph{optimistic} guess for real scenarios (finite burns always cost a little more).
\end{itemize}

\paragraph{Effect of an impulse.} At $P$ (true anomaly $\theta_*^-$ on the old orbit):
\begin{equation}
\underline r^+=\underline r^-,\qquad \underline v^+=\underline v^-+\Delta\underline v,\qquad
\Delta\underline v=\Delta v_\theta\,\hat\theta+\Delta v_r\,\hat r .
\end{equation}
With $v_r^-=\sqrt{\mu/p^-}\,e^-s_{\theta_*^-}$ and $v_\theta^-=\sqrt{\mu/p^-}\,(1+e^-c_{\theta_*^-})$:
\begin{equation}
v_r^+=v_r^-+\Delta v_r,\qquad v_\theta^+=v_\theta^-+\Delta v_\theta .
\end{equation}
The new orbit follows in three steps:
\begin{enumerate}
\item $a^+$ from the \textbf{energy}: $\dfrac{(v^+)^2}{2}-\dfrac{\mu}{r}=-\dfrac{\mu}{2a^+}\Rightarrow a^+=\dfrac{\mu}{2\mu/r-(v^+)^2}$;
\item $e^+$ from the \textbf{angular momentum}: $h^+=r\,v_\theta^+=\sqrt{\mu a^+[1-(e^+)^2]}\Rightarrow e^+=\sqrt{1-\dfrac{h^{+2}}{\mu a^+}}$;
\item $\theta_*^+$ from \textbf{radius and radial velocity}: $c_{\theta_*^+}=\dfrac{p^+-r}{r\,e^+}$, $s_{\theta_*^+}=\dfrac{v_r^+}{e^+}\sqrt{\dfrac{p^+}{\mu}}$
$\Rightarrow\theta_*^+=2\operatorname{atan}\dfrac{s_{\theta_*^+}}{1+c_{\theta_*^+}}$ (correct quadrant).
\end{enumerate}
The apsidal line rotates by $\theta_*^--\theta_*^+$; the vacant focus moves so that the new velocity satisfies the reflection
property of the new ellipse. A tangential impulse ($\Delta v_r=0$ at an apse) does not rotate the apsidal line.

\begin{center}
\includegraphics[width=0.55\linewidth]{impulse}\\
{\footnotesize Same position, new velocity: the orbit through $P$ changes shape and orientation.}
\end{center}

\paragraph{Examples.}
\begin{itemize}
\item \textbf{Valid}: orbit transfers with chemical engines (e.g.\ LEO$\to$GEO Hohmann transfer: burns of a few minutes vs
orbital periods of hours; plane changes at the nodes; injection into an escape hyperbola; rendezvous burns). The
Bounded Multi Impulse model, used to approximate finite burns, is itself based on impulsive $\Delta v$'s.
\item \textbf{Not valid}: (a) the ascent trajectory of a launch vehicle (thrust acts for minutes while gravity and drag
cause large losses and the position changes by hundreds of km); (b) low-thrust trajectories (electric propulsion), where
continuous thrust is applied for long time intervals (weeks or months) and the orbit spirals.
\end{itemize}

\begin{keyf}
$T\Delta t=\Delta v$ ($T\to\infty$, $\Delta t\to0$), \quad $\underline r^+=\underline r^-$, \quad $\underline v^+=\underline v^-+\Delta\underline v$,
\quad $a^+=\dfrac{\mu}{2\mu/r-(v^+)^2}$, \quad $e^+=\sqrt{1-\dfrac{(r\,v_\theta^+)^2}{\mu a^+}}$.
\end{keyf}

\pitfalls
\begin{itemize}
\item The position does \emph{not} change: the old and new orbits intersect at $P$.
\item Say that the approximation is optimistic and requires a burn time much shorter than the orbital period.
\item Use the radial \emph{and} horizontal components; $\theta_*^+$ needs both $\sin$ and $\cos$ for the quadrant.
\end{itemize}

\source{cap03.tex l.18--43 (approximation, validity), l.45--113 (effect of velocity impulses). Formulario: \fml{3-IMP}, \fml{3-NEW}.}

% ---------------------------------------------------------------------
\question{A6}{Hohmann transfer vs bielliptic transfer}
% source: cap03.tex l.115-377
\begin{qbox}
Compare a Hohmann transfer with a bielliptic transfer (between circular coplanar orbits), discussing the conditions under
which each option is optimal.
\qmeta{$\times$2 in the list (2025); Jan 2025; Jun 2025 ("discussing the conditions under which each option is optimal").}
\end{qbox}

\begin{outline}
\begin{enumerate}
\item Problem: coplanar circular orbits $R_i\to R_f$, minimise $\Delta v_{tot}$ (= propellant).
\item Hohmann: bitangent ellipse, $a_T$, $e_T$, $\Delta v_1$, $\Delta v_2$, $\Delta t=\pi\sqrt{a_T^3/\mu}$; globally optimal among 2-impulse transfers (sketch of proof).
\item Bielliptic: 3 impulses, intermediate apoapsis $r_A^{\max}$, $\Delta v_{1,2,3}$, $\Delta t$; biparabolic limit.
\item Non-dimensional comparison $\sigma=R_f/R_i$, $\rho=r_A^{\max}/R_i\ge\sigma$; plot.
\item Result: $\sigma<11.94$ Hohmann; $11.94<\sigma<15.58$ depends on $\rho$; $\sigma>15.58$ bielliptic. Trade-off with time of flight.
\end{enumerate}
\end{outline}

\answer
\paragraph{Problem.} Given two coplanar circular orbits of radii $R_i$ and $R_f$ ($R_f>R_i$), find the impulsive transfer that
minimises $\Delta v_{tot}$, i.e.\ the propellant consumption (Tsiolkovsky).

\paragraph{Hohmann transfer.} The transfer arc is an ellipse tangent to both circles (bitangent), with periapsis on the
initial orbit and apoapsis on the final one:
\begin{equation}
a_{T}=\frac{R_{i}+R_{f}}{2},\qquad e_{T}=\frac{R_{f}-R_{i}}{R_{f}+R_{i}} .
\end{equation}
Two impulses, both \emph{tangent and in the direction of motion} (both increase $a$):
\begin{equation}
\Delta v_{1}=\sqrt{\frac{2\mu}{R_{i}+R_{f}}\frac{R_{f}}{R_{i}}}-\sqrt{\frac{\mu}{R_{i}}},\qquad
\Delta v_{2}=\sqrt{\frac{\mu}{R_{f}}}-\sqrt{\frac{2\mu}{R_{i}+R_{f}}\frac{R_{i}}{R_{f}}},\qquad
\Delta t=\pi\sqrt{\frac{a_{T}^{3}}{\mu}} .
\end{equation}
\emph{Optimality.} Among all two-impulse transfers with an elliptic arc, the feasible ones satisfy $r_{P_T}\le R_i$ and
$r_{A_T}\ge R_f$. For each impulse
$(\Delta v_k)^2=\frac{3\mu}{R_k}-2\frac{\mu}{R_k}\sqrt{\frac{p_T}{R_k}}-\frac{\mu}{p_T}(1-e_T^2)$, and
$\partial(\Delta v_k)^2/\partial e_T=2e_T\mu/p_T>0$: at fixed $p_T$ the cost grows with $e_T$; along both boundaries of the
feasible region $d(\Delta v_1+\Delta v_2)/de_T>0$ too. Hence the minimum is at the corner point $H$ where both constraints
are active, the bitangent ellipse: the Hohmann transfer is the \textbf{globally optimal two-impulse} transfer.

With $\sigma=R_f/R_i$:
$\dfrac{\Delta v^{(H)}_{tot}}{\sqrt{\mu/R_i}}=\sqrt{\dfrac{2\sigma}{1+\sigma}}-1+\dfrac{1}{\sqrt\sigma}-\sqrt{\dfrac{2}{\sigma(1+\sigma)}}$;
it has a \textbf{maximum at $\sigma\simeq15.58$} and the asymptote $\sqrt2-1$ (transfer to a parabola) as $\sigma\to\infty$.

\paragraph{Bielliptic transfer.} Three impulses and two half-ellipses with a common apoapsis $r_A^{\max}\ge R_f$:
\begin{equation}
\sys{\Delta v_{1}^{(BE)}&=\sqrt{\tfrac{2\mu}{R_{i}+r_{A}^{\max}}\tfrac{r_{A}^{\max}}{R_{i}}}-\sqrt{\tfrac{\mu}{R_{i}}}
&&\text{(at $R_i$, along the motion)}\\
\Delta v_{2}^{(BE)}&=\sqrt{\tfrac{2\mu}{R_{f}+r_{A}^{\max}}\tfrac{R_{f}}{r_{A}^{\max}}}-\sqrt{\tfrac{2\mu}{R_{i}+r_{A}^{\max}}\tfrac{R_{i}}{r_{A}^{\max}}}
&&\text{(at $r_A^{\max}$: raises the periapsis to $R_f$)}\\
\Delta v_{3}^{(BE)}&=\sqrt{\tfrac{2\mu}{R_{f}+r_{A}^{\max}}\tfrac{r_{A}^{\max}}{R_{f}}}-\sqrt{\tfrac{\mu}{R_{f}}}
&&\text{(at $R_f$, \emph{opposite} to the motion: circularises)}}
\end{equation}
\begin{equation}
\Delta t=\pi\sqrt{\frac{a_{T1}^{3}}{\mu}}+\pi\sqrt{\frac{a_{T2}^{3}}{\mu}},\qquad a_{T1}=\frac{R_i+r_A^{\max}}{2},\ a_{T2}=\frac{R_f+r_A^{\max}}{2}.
\end{equation}
As $r_A^{\max}$ increases $\Delta v_{tot}$ decreases but $\Delta t$ increases. The limit $r_A^{\max}\to\infty$ is the
\textbf{biparabolic transfer}: $\Delta v_1=(\sqrt2-1)\sqrt{\mu/R_i}$, $\Delta v_2=0$, $\Delta v_3=(\sqrt2-1)\sqrt{\mu/R_f}$,
$\Delta t=\infty$ (best of all bielliptic transfers, but infeasible). The idea: the expensive part of the plane-raising is done
far away, where the velocity is small and a small $\Delta v_2$ raises the periapsis a lot.

\begin{center}
\includegraphics[width=0.9\linewidth]{hoh_be_geom}
\end{center}

\paragraph{Comparison.} With $\sigma=R_f/R_i$ and $\rho=r_A^{\max}/R_i\ge\sigma$ both costs are functions of $(\sigma,\rho)$ only:
\begin{center}
\includegraphics[width=0.72\linewidth]{hoh_be_dv}\\
{\footnotesize $\Delta v_{tot}/\sqrt{\mu/R_i}$ computed from the formulas above; bielliptic curves exist only for $\sigma\le\rho$.}
\end{center}
\begin{itemize}
\item Hohmann and biparabolic curves intersect at $\sigma=11.94$; both tend to $\sqrt2-1$.
\item For finite $\rho$ the intersection with the Hohmann curve moves towards $\sigma=15.58$ as $\rho$ decreases: all
intersections lie in $11.94\le\sigma\le15.58$.
\end{itemize}
\begin{equation}
\begin{cases}
\sigma<11.94: & \text{Hohmann is globally optimal;}\\
11.94<\sigma<15.58: & \text{either may be optimal, depending on }\rho=r_A^{\max}/R_i\ge\sigma;\\
\sigma>15.58: & \text{the bielliptic transfer is optimal (for any }\rho\ge\sigma).
\end{cases}
\end{equation}

\begin{center}
\small
\begin{tabular}{@{}p{3.4cm}p{5.6cm}p{5.6cm}@{}}
\toprule
 & Hohmann & Bielliptic\\
\midrule
impulses & 2, both along the motion & 3 (two along, the last against the motion)\\
free parameter & none & $r_A^{\max}$ (limited by the mission time)\\
$\Delta v_{tot}$ & minimum for $\sigma<11.94$ & lower for $\sigma>15.58$; gain small (few \%)\\
time of flight & half period of one ellipse & two half periods, much longer (grows with $r_A^{\max}$)\\
optimality & globally optimal among 2-impulse transfers & beats Hohmann only for large $\sigma$\\
practical use & almost always (e.g.\ LEO$\to$GEO: $\sigma\approx6.3$) & rare: large $\sigma$ and no time constraint\\
\bottomrule
\end{tabular}
\end{center}

\begin{keyf}
$a_T=\dfrac{R_i+R_f}{2}$, \ $\Delta v_1=\sqrt{\dfrac{2\mu R_f}{R_i(R_i+R_f)}}-\sqrt{\dfrac{\mu}{R_i}}$, \
$\Delta v_2=\sqrt{\dfrac{\mu}{R_f}}-\sqrt{\dfrac{2\mu R_i}{R_f(R_i+R_f)}}$, \ $\Delta t=\pi\sqrt{a_T^3/\mu}$;
thresholds $\sigma=11.94$ and $15.58$; biparabolic $\Delta v=(\sqrt2-1)(\sqrt{\mu/R_i}+\sqrt{\mu/R_f})$.
\end{keyf}

\pitfalls
\begin{itemize}
\item Hohmann is optimal among \emph{two}-impulse transfers only; three-impulse transfers can do better for large $\sigma$.
\item $\Delta v_3$ of the bielliptic transfer is a braking impulse (opposite to the motion).
\item Always mention the time of flight: the gain of the bielliptic transfer is paid with a much longer transfer.
\item The 15.58 value is also the maximum of the Hohmann cost; 11.94 is the Hohmann--biparabolic crossing.
\end{itemize}

\source{cap03.tex l.115--240 (two-impulse optimality), l.242--284 (Hohmann), l.286--344 (bielliptic, biparabolic), l.346--377
(comparison). Thresholds re-checked numerically in \texttt{figures/src/make\_figures.py}. Formulario: \fml{3-HOH}, \fml{3-BEL}, \fml{3-HVB}.}

% ---------------------------------------------------------------------
\question{A7}{Globally optimal transfers from an elliptic orbit to a hyperbolic path}
% source: cap03.tex l.379-449
\begin{qbox}
Describe two optimal ways to transfer a spacecraft from an elliptical orbit to a hyperbolic trajectory and discuss their
global optimality.
\qmeta{list by chapter (Ch.\,3).}
\end{qbox}

\begin{outline}
\begin{enumerate}
\item Problem: planar; initial ellipse ($r_p$, $r_A$); final hyperbola/parabola with given energy $\mathcal E_f\ge0$; no orientation constraint;
bounds $r_A^{\max}$, $r_p^{\min}$ on intermediate orbits.
\item Why periapsis: at fixed $r$, $\Delta\mathcal E=v\,\Delta v+\Delta v^2/2$ $\Rightarrow$ impulses are most effective where $v$ is largest.
\item Case 1, $\mathcal E_f<\mu/r_A^{\max}$: one tangential impulse at periapsis.
\item Case 2, $\mathcal E_f>\mu/r_A^{\max}$: three impulses (raise apoapsis to $r_A^{\max}$, lower periapsis to $r_p^{\min}$, escape at the new periapsis).
\item Global optimality: these are the globally optimal solutions of the problem; analogy with Hohmann/bielliptic; application to interplanetary injection.
\end{enumerate}
\end{outline}

\answer
\paragraph{Problem statement.} Planar transfer, with no constraint on the orientation of the terminal orbits:
\begin{itemize}
\item initial ellipse with specified $r_A$ and $r_p$;
\item final hyperbola ($\mathcal E_f>0$) or parabola ($\mathcal E_f=0$) with specified specific energy $\mathcal E_f$,
coplanar with the ellipse (only the energy, i.e.\ $v_\infty=\sqrt{2\mathcal E_f}$, matters, not the direction);
\item intermediate orbits are bounded: apoapsis radius $\le r_A^{\max}$ (e.g.\ sphere of influence, time) and periapsis radius
$\ge r_p^{\min}$ (e.g.\ planetary radius plus atmosphere).
\end{itemize}

\paragraph{Why impulses at periapsis.} An impulse does not change $r$, so from $\mathcal E=v^2/2-\mu/r$ a tangential impulse gives
\begin{equation}
\Delta\mathcal E=\frac{(v+\Delta v)^2-v^2}{2}=v\,\Delta v+\frac{\Delta v^2}{2}:
\end{equation}
the same $\Delta v$ produces the largest energy increase where the speed is largest, i.e.\ at periapsis; and the lower the
periapsis, the larger the speed there.

\paragraph{The two globally optimal strategies.}
\begin{enumerate}[label=(\arabic*)]
\item $\mathcal E_f<\mu/r_A^{\max}$: the globally optimal transfer is a \textbf{single tangential impulse at periapsis} that raises
the energy to $\mathcal E_f$:
\begin{equation}
\Delta v=\sqrt{2\left(\mathcal E_f+\frac{\mu}{r_p}\right)}-v_p,\qquad v_p=\sqrt{\frac{\mu}{a}\frac{1+e}{1-e}} .
\end{equation}
\item $\mathcal E_f>\mu/r_A^{\max}$: the globally optimal transfer has \textbf{three impulses}:
\begin{enumerate}[label=(\alph*)]
\item impulse at periapsis, to raise the apoapsis radius to $r_A^{\max}$;
\item impulse at apoapsis (where the speed is small, so it is cheap), to lower the periapsis radius to $r_p^{\min}$;
\item impulse at the new periapsis $r_p^{\min}$, where the speed is now much higher, to increase the energy to $\mathcal E_f$.
\end{enumerate}
\end{enumerate}
\begin{center}
\includegraphics[width=0.92\linewidth]{ell_hyp}
\end{center}

\paragraph{Global optimality (discussion).} The two solutions are the global optima of the problem above, i.e.\ no other
sequence of impulses within the bounds yields a smaller $\Delta v_{tot}$; which one applies depends only on the comparison
between $\mathcal E_f$ and $\mu/r_A^{\max}$. The three-impulse strategy pays two extra impulses to "move" the final burn
to a lower periapsis, which is convenient only when the final energy (the escape burn) is large enough; otherwise the single
periapsis burn is optimal. The logic is the same as for circle-to-circle transfers: the Hohmann transfer is optimal for
moderate changes, the bielliptic one (going far away, where velocity changes are cheap) for large changes. The limiting
values $r_A^{\max}$, $r_p^{\min}$ come from practical constraints (time, sphere of influence, atmosphere, safety).

\paragraph{Application: injection into an interplanetary hyperbola.} From a circular parking orbit of radius $r_p$
the required hyperbolic excess speed $v_\infty$ (from the heliocentric transfer, e.g.\ $v_\infty^{(E)}=\Delta v_1^{(H)}$) fixes
$\mathcal E_f=v_\infty^2/2$, and the single periapsis burn is
$\Delta v=\sqrt{v_\infty^2+2\mu/r_p}-\sqrt{\mu/r_p}$, which is smaller for lower parking orbits.

\begin{keyf}
$\Delta\mathcal E=v\Delta v+\tfrac12\Delta v^2$; \ case 1 ($\mathcal E_f<\mu/r_A^{\max}$): $\Delta v=\sqrt{2(\mathcal E_f+\mu/r_p)}-v_p$;
\ case 2: periapsis $\to r_A^{\max}$, apoapsis $\to r_p^{\min}$, periapsis $\to\mathcal E_f$; \ $v_p=\sqrt{v_\infty^2+2\mu/r_p}$.
\end{keyf}

\pitfalls
\begin{itemize}
\item Specify the assumptions: planar, no constraint on the orientation of the hyperbola, only $\mathcal E_f$ prescribed.
\item State the threshold $\mu/r_A^{\max}$ that separates the two cases.
\item Explain physically why burns are done at periapsis (energy gain proportional to the speed).
\end{itemize}

\source{cap03.tex l.379--401 (globally optimal transfer from ellipse to hyperbola), l.403--449 (injection into a hyperbolic path).
The energy-gain argument $\Delta\mathcal E=v\Delta v+\Delta v^2/2$ follows from the energy integral. Formulario: \fml{3-EHY}, \fml{3-INJ}.}

% ---------------------------------------------------------------------
\question{A8}{Losses with respect to Tsiolkovsky: the $\Delta V$-loss integrals}
% source: cap06.tex l.112-162 (Tsiolkovsky), l.504-571 (losses equation)
\begin{qbox}
Losses with respect to the Tsiolkovsky equation: define the $\Delta V$-loss integrals and explain their physical meaning.
\qmeta{$\times$2 in the list; Sep 2024.}
\qnote{Difference from \qref{A9}: here the focus is on the derivation and meaning of the integrals; A9 adds the
propulsive velocity and how each loss depends on the trajectory and on the vehicle.}
\end{qbox}

\begin{outline}
\begin{enumerate}
\item Tsiolkovsky: assumptions (no gravity, no drag, thrust along $\underline v$) $\Rightarrow$ ideal upper bound $\Delta v=c\ln(m_0/m_f)$.
\item General equation $\dot{\underline v}=(\underline G+\underline A+\underline T)/m$, projection on $\hat v$ with $\gamma$, $\alpha$, $D$ (sketch).
\item Rearrange and integrate: $\Delta v=\Delta v_T-\Delta v_G-\Delta v_M-\Delta v_D$ (the four integrals).
\item Physical meaning of each term; typical values (VEGA); when each one dominates.
\end{enumerate}
\end{outline}

\answer
\paragraph{Tsiolkovsky's law and its assumptions.} With (i) no aerodynamic force, (ii) no gravity, (iii) thrust aligned with
the velocity: $\dot v=T/m$, $\dot m=-T/c$ $\Rightarrow$ $\dot v=-c\,\dot m/m$, whose integral is
\begin{equation}
\Delta v=v_f-v_0=c\ln\frac{m_0}{m_f}\qquad(c=g_0I_{sp}).
\end{equation}
It gives the \emph{maximum} $\Delta v$ obtainable from a given mass variation (the limiting performance). For a launch vehicle
the assumptions are violated, and the actual velocity gain is smaller: the differences are the \emph{losses}.

\paragraph{Losses equation.} The general vector equation is
\begin{equation}
\frac{d\underline v}{dt}=\frac{\underline G+\underline A+\underline T}{m},\qquad \underline G=-\frac{\mu}{r^2}\hat r\ (\text{per unit mass: } m\underline g) .
\end{equation}
Project it on $\hat v$. Let $\gamma$ = flight-path angle (between $\underline v$ and the local horizontal), $\alpha$ = misalignment
(incidence) angle between $\underline T$ and $\underline v$, $D$ = drag magnitude (lift and side force neglected):
\begin{center}\includegraphics[width=0.36\linewidth]{losses}\end{center}
\begin{equation}
\frac{dv}{dt}=-\frac{\mu}{r^{2}}\,s_{\gamma}+\frac{T}{m}\,c_{\alpha}-\frac{D}{m}
=\frac{T}{m}-\frac{\mu}{r^{2}}\,s_{\gamma}-\frac{T}{m}\left(1-c_{\alpha}\right)-\frac{D}{m}.
\end{equation}
Integrating from $t_0$ to $t_f$:
\begin{equation}
\Delta v=\underbrace{\int_{t_0}^{t_f}\!\frac{T}{m}\,dt}_{\Delta v_T}
-\underbrace{\int_{t_0}^{t_f}\!\frac{\mu}{r^{2}}\,s_{\gamma}\,dt}_{\Delta v_G}
-\underbrace{\int_{t_0}^{t_f}\!\frac{T}{m}\left(1-c_{\alpha}\right)dt}_{\Delta v_M}
-\underbrace{\int_{t_0}^{t_f}\!\frac{D}{m}\,dt}_{\Delta v_D}.
\end{equation}

\paragraph{Physical meaning.}
\begin{itemize}
\item $\Delta v_T$ -- \textbf{ideal (Tsiolkovsky) velocity change}: since $T=-c\,\dot m$, $\Delta v_T=c\ln(m_0/m_f)$. It depends only
on the engine ($c$) and on the mass ratio, not on the trajectory.
\item $\Delta v_G$ -- \textbf{gravity loss}: the component of gravity along the velocity, $-g\,s_\gamma$, decelerates the vehicle
while it climbs (thrust spent to hold the vehicle up against gravity instead of accelerating it). It is large when the path is
steep ($\gamma$ near $90^\circ$) and the burn is long; it vanishes for horizontal flight.
\item $\Delta v_M$ -- \textbf{misalignment loss}: only the component $T c_\alpha$ changes the speed; the normal component
$T s_\alpha$ only rotates the velocity. It vanishes for $\alpha=0$ (gravity turn, thrust along $\underline v$).
\item $\Delta v_D$ -- \textbf{drag loss}: the work of the aerodynamic drag, which dissipates kinetic energy; it depends on
density, speed and on $S/m$ ($D/m=\frac12\rho v_R^2C_DS/m$).
\end{itemize}
All losses are positive integrals, so $\Delta v<\Delta v_T$. The equation holds for each sub-rocket.

\paragraph{Typical values and trends.} For VEGA: $\Delta v_T=3.054$, $\Delta v_G=0.683$, $\Delta v_M=0.019$, $\Delta v_D=0.107$\,km/s:
\textbf{gravity loss dominates}. In the early ascent $\gamma$ is large and $\Delta v_G$ dominates; a gravity-turn condition
($\alpha=0$) gives $\Delta v_M=0$; as velocity grows $D$ grows while $\gamma$ decreases, reducing $\Delta v_G$. Hence:
(a) the thick atmospheric layers are crossed along a \emph{near-radial} path, to limit drag and minimise the time spent in
them; (b) as soon as the density decreases the velocity is rotated to reduce $\gamma$, to avoid excessive gravity losses.

\begin{keyf}
$\dfrac{dv}{dt}=-\dfrac{\mu}{r^2}s_\gamma+\dfrac{T}{m}c_\alpha-\dfrac{D}{m}$; \
$\Delta v=\Delta v_T-\Delta v_G-\Delta v_M-\Delta v_D$; \ $\Delta v_T=\displaystyle\int\frac{T}{m}dt=c\ln\frac{m_0}{m_f}$, \
$\Delta v_G=\displaystyle\int\frac{\mu}{r^2}s_\gamma dt$, \ $\Delta v_M=\displaystyle\int\frac{T}{m}(1-c_\alpha)dt$, \ $\Delta v_D=\displaystyle\int\frac{D}{m}dt$.
\end{keyf}

\pitfalls
\begin{itemize}
\item Define $\gamma$ and $\alpha$ with a sketch; the gravity term has $s_\gamma$ (not $c_\gamma$) because $\gamma$ is measured from the horizontal.
\item The misalignment loss is $T(1-c_\alpha)/m$, not $T s_\alpha/m$.
\item Recall that $\Delta v_T$ is exactly Tsiolkovsky's $\Delta v$ and that the losses are always subtracted.
\end{itemize}

\source{cap06.tex l.112--162 (Tsiolkovsky's law), l.504--571 (losses equation, VEGA values, comments). Formulario: \fml{6-TSI}.}

% ---------------------------------------------------------------------
\question{A9}{Propulsive velocity and velocity losses in the ascent trajectory}
% source: cap06.tex l.356-502 (ascent phases), l.504-571 (losses)
\begin{qbox}
Define propulsive velocity and the velocity losses due to misalignment, gravity and aerodynamic effects in the ascent
trajectory of the launch vehicle. Explain how each type of velocity loss depends on vehicle trajectory and possibly other
launch vehicle features.
\qmeta{Jul 2025 (listed as a 2025 question).}
\qnote{Difference from \qref{A8}: same integrals; here the emphasis is on the dependence on the trajectory and the vehicle.}
\end{qbox}

\begin{outline}
\begin{enumerate}
\item Equation along $\hat v$ and the four integrals (sketch with $\gamma$, $\alpha$).
\item Propulsive (ideal) velocity $\Delta v_T=\int T/m\,dt=c\ln(m_0/m_f)$: engine and mass ratio only.
\item Gravity loss: depends on $\gamma(t)$ and burn time $\Rightarrow$ thrust-to-weight ratio, steepness of the path.
\item Misalignment loss: depends on $\alpha(t)$ $\Rightarrow$ gravity turn, structural limit $q\alpha$.
\item Drag loss: depends on $\rho(h)$, $v$, ballistic coefficient $C_DS/m$ $\Rightarrow$ fast near-vertical crossing.
\item Trade-off steep vs shallow ascent; phases of the ascent; VEGA numbers.
\end{enumerate}
\end{outline}

\answer
\paragraph{Losses equation.} Projecting $\dot{\underline v}=(\underline G+\underline A+\underline T)/m$ on the velocity direction
($\gamma$ = flight-path angle from the local horizontal, $\alpha$ = angle between thrust and velocity, $D$ = drag):
\begin{equation}
\frac{dv}{dt}=\frac{T}{m}-\frac{\mu}{r^{2}}\,s_{\gamma}-\frac{T}{m}\left(1-c_{\alpha}\right)-\frac{D}{m}
\;\Rightarrow\;
\Delta v=\Delta v_T-\Delta v_G-\Delta v_M-\Delta v_D .
\end{equation}
\begin{center}\includegraphics[width=0.33\linewidth]{losses}\end{center}

\paragraph{Propulsive velocity.} The propulsive (ideal) velocity is the velocity change that the thrust alone would produce:
\begin{equation}
\Delta v_T=\int_{t_0}^{t_f}\frac{T}{m}\,dt=\int_{t_0}^{t_f}-c\frac{\dot m}{m}\,dt=c\ln\frac{m_0}{m_f}\qquad(c=g_0I_{sp}).
\end{equation}
It coincides with Tsiolkovsky's $\Delta v$; it depends only on the engine (effective exhaust velocity) and on the mass
ratio of the stage, \emph{not} on the trajectory. The actual gain is $\Delta v_T$ minus the three losses.

\paragraph{Gravity loss} $\Delta v_G=\displaystyle\int\frac{\mu}{r^2}s_\gamma\,dt$.
\begin{itemize}
\item \emph{Trajectory}: proportional to $s_\gamma$: maximum during the vertical arc ($\gamma=90^\circ$), small when the path is
nearly horizontal. Early pitch-over / lower $\gamma$ reduce it.
\item \emph{Vehicle}: it is an integral over time, so it grows with the \emph{burn duration}: a high thrust-to-weight ratio
(short burns) reduces it, a low-thrust upper stage that burns long while climbing increases it.
\item Usually the largest loss (VEGA: 0.683\,km/s out of 3.054).
\end{itemize}

\paragraph{Misalignment loss} $\Delta v_M=\displaystyle\int\frac{T}{m}(1-c_\alpha)\,dt$.
\begin{itemize}
\item \emph{Trajectory/steering}: zero when the thrust is aligned with the velocity ($\alpha=0$), as in the \textbf{gravity turn}, where
gravity alone turns the velocity. It appears in the pitch manoeuvre and in the rarefied-atmosphere arcs where an incidence is used
to shape the trajectory.
\item \emph{Vehicle}: the admissible $\alpha$ is limited by structural loads, $q\alpha\le(q\alpha)_{\max}$ with $q=\frac12\rho v_R^2$; vehicles
that tolerate an incidence (military-derived) can steer more but pay this loss. It is the smallest (VEGA: 0.019\,km/s).
\end{itemize}

\paragraph{Drag loss} $\Delta v_D=\displaystyle\int\frac{D}{m}\,dt$, \quad $\dfrac{D}{m}=\dfrac12\rho v_R^2\,C_D\dfrac{S}{m}$.
\begin{itemize}
\item \emph{Trajectory}: depends on the density along the path ($\rho$ decreases exponentially with altitude) and on the speed:
the dense layers (below $\sim40$\,km) must be crossed \emph{quickly and near-vertically}, before reaching high speeds.
\item \emph{Vehicle}: proportional to the ballistic coefficient $C_DS/m$ (shape, cross-section, mass); large, heavy vehicles have
smaller $S/m$ and hence smaller drag losses \notinnotes. VEGA: 0.107\,km/s.
\end{itemize}

\paragraph{Trade-off and ascent phases.} A steep path reduces drag losses but increases gravity losses; a shallow path does the
opposite. Real ascents: vertical arc at launch (launch-site safety; avoid premature rotation of $\underline v$), roll and pitch
manoeuvres (select the orbit plane; some misalignment), \textbf{gravity turn} ($\alpha=0$ to avoid destructive bending loads:
$\Delta v_M\approx0$), rarefied-atmosphere arc (small $\alpha$ allowed), coast arc, injection arc (thrust nearly horizontal:
$\Delta v_G\approx0$). Hence: near-radial crossing of the thick atmosphere, then rotation of $\underline v$ to reduce $\gamma$.

\begin{keyf}
$\Delta v_T=\displaystyle\int\frac Tm dt=c\ln\frac{m_0}{m_f}$; \ $\Delta v_G=\displaystyle\int g\,s_\gamma dt$ (steepness, burn time);
\ $\Delta v_M=\displaystyle\int\frac Tm(1-c_\alpha)dt$ (steering, $q\alpha$ limit); \ $\Delta v_D=\displaystyle\int\frac12\rho v^2C_D\frac{S}{m}dt$ (density, speed, $C_DS/m$).
\end{keyf}

\pitfalls
\begin{itemize}
\item "Propulsive velocity" is the ideal $\Delta v_T$, not the final orbital speed.
\item For each loss say both the trajectory dependence and a vehicle feature (thrust-to-weight, structural $q\alpha$ limit, $C_DS/m$).
\item Mention the conflicting requirements (drag vs gravity) and the gravity turn.
\end{itemize}

\source{cap06.tex l.356--482 (phases of the ascent), l.504--571 (losses). Formulario: \fml{6-TSI}.}

% ---------------------------------------------------------------------
\question{A10}{Rocket staging: mass-distribution parameters and motivation}
% source: cap06.tex l.164-224, 225-354
\begin{qbox}
With regard to rocket staging, define the key mass-distribution parameters of each stage and briefly justify the use of
staging for launch vehicles.
\qmeta{list by chapter (Ch.\,6).}
\qnote{Adjacent to \qref{A11} (same parameters; A11 asks for the total $\Delta V$ formula).}
\end{qbox}

\begin{outline}
\begin{enumerate}
\item Sub-rocket $i$: $m_0^{(i)}=m_s^{(i)}+m_p^{(i)}+m_u^{(i)}$, $m_u^{(i)}=m_0^{(i+1)}$ (sketch).
\item Structural coefficient $\epsilon_s^{(i)}$ (technology, $\ge\epsilon_{s,\min}$) and mass ratio $u_0^{(i)}=m_f^{(i)}/m_0^{(i)}$.
\item Masses in terms of $(\epsilon_s,u_0)$; $\Delta v^{(i)}=-c_i\ln u_0^{(i)}$; $m_u>0$ iff $u_0>\epsilon_s$; low $\epsilon_s$ desirable.
\item Why staging: discard useless structural mass; single-stage limit $c\ln(1/\epsilon_s)$; $\Delta v_{\max}(N)$ with diminishing returns, 3--4 stages.
\end{enumerate}
\end{outline}

\answer
\paragraph{Definitions.} A multistage launch vehicle is a sequence of sub-rockets. Sub-rocket $i$ has initial mass
\begin{equation}
m_{0}^{(i)}=m_{s}^{(i)}+m_{p}^{(i)}+m_{u}^{(i)},
\end{equation}
$m_s^{(i)}$ = structural (inert) mass of stage $i$ (tanks, engines, structure), $m_p^{(i)}$ = propellant mass,
$m_u^{(i)}$ = "payload" of sub-rocket $i$, i.e.\ the initial mass of the next sub-rocket, $m_u^{(i)}=m_0^{(i+1)}$, and the true
payload for the last one.
\begin{center}\includegraphics[width=0.7\linewidth]{staging}\end{center}
The mass distribution of each stage is described by two parameters:
\begin{equation}
\epsilon_{s}^{(i)}=\frac{m_{s}^{(i)}}{m_{s}^{(i)}+m_{p}^{(i)}}\quad\text{(structural coefficient)},\qquad
u_{0}^{(i)}=\frac{m_{f}^{(i)}}{m_{0}^{(i)}}=\frac{m_{s}^{(i)}+m_{u}^{(i)}}{m_{0}^{(i)}}\quad\text{(mass ratio)}.
\end{equation}
$\epsilon_s$ depends on the \textbf{technology} (materials, tank pressure, engine mass), so $\epsilon_s^{(i)}\ge\epsilon_{s,\min}^{(i)}$;
$u_0$ is the final-to-initial mass ratio of the burn. Inverting the definitions:
\begin{equation}
m_{p}^{(i)}=m_{0}^{(i)}\left[1-u_{0}^{(i)}\right],\quad
m_{s}^{(i)}=m_{0}^{(i)}\frac{\epsilon_{s}^{(i)}\left[1-u_{0}^{(i)}\right]}{1-\epsilon_{s}^{(i)}},\quad
m_{u}^{(i)}=m_{0}^{(i)}\frac{u_{0}^{(i)}-\epsilon_{s}^{(i)}}{1-\epsilon_{s}^{(i)}},
\end{equation}
and the velocity change of stage $i$ (Tsiolkovsky) is $\Delta v^{(i)}=-c_i\ln u_0^{(i)}$.
\begin{itemize}
\item $m_u^{(i)}>0$ only if $u_0^{(i)}>\epsilon_s^{(i)}$; for $u_0=\epsilon_s$ the stage carries nothing.
\item At given $u_0$ (i.e.\ given $\Delta v^{(i)}$) the payload fraction $m_u/m_0$ is maximum for $\epsilon_s=\epsilon_{s,\min}$:
\textbf{low structural coefficients} improve performance.
\end{itemize}

\paragraph{Why staging.} The purpose is to increase the overall $\Delta v$ by \textbf{discarding structural masses that have become
useless} (empty tanks and their engines). With a single stage the whole structure must be accelerated up to the final speed;
even with zero payload ($u_0\to\epsilon_s$) the maximum is
\begin{equation}
\Delta v_{\max}^{\text{single}}=c\ln\frac{1}{\epsilon_s},
\end{equation}
e.g.\ $c=3$\,km/s, $\epsilon_s=0.1$ $\Rightarrow$ 6.9\,km/s, below the $\approx9$--10\,km/s needed to reach LEO including
losses \notinnotes. With $N$ stages the velocity changes add up, $\Delta v=-\sum_i c_i\ln u_0^{(i)}$; for similar stages the optimal
distribution (all $u_0^{(i)}$ equal) gives
\begin{equation}
\Delta v_{\max}=-c\,N\ln\left\{\epsilon_{s}+(1-\epsilon_{s})\left[\frac{m_{u}^{(N)}}{m_{0}^{(1)}}\right]^{1/N}\right\},
\end{equation}
which increases with $N$ with diminishing returns and has a finite limit for $N\to\infty$. Structural and operational complexity
grow with $N$: launch vehicles rarely have more than 5 stages (often 3--4). Staging also allows engines optimised for
different altitudes (sea-level vs vacuum nozzles) \notinnotes.

\begin{keyf}
$m_0^{(i)}=m_s^{(i)}+m_p^{(i)}+m_u^{(i)}$, \ $\epsilon_s=\dfrac{m_s}{m_s+m_p}$, \ $u_0=\dfrac{m_s+m_u}{m_0}$, \
$\dfrac{m_u}{m_0}=\dfrac{u_0-\epsilon_s}{1-\epsilon_s}$, \ $\Delta v^{(i)}=-c_i\ln u_0^{(i)}$.
\end{keyf}

\pitfalls
\begin{itemize}
\item $\epsilon_s$ uses structure + propellant only (no payload); $u_0$ uses the final mass, which includes the upper stages.
\item The "payload" of stage $i$ is the whole upper composite, $m_u^{(i)}=m_0^{(i+1)}$.
\item Justify staging with the single-stage limit and the diminishing returns with $N$.
\end{itemize}

\source{cap06.tex l.164--224 (rocket staging), l.225--354 (optimal staging, plot of $\Delta v_{\max}$ vs $N$; source errata \#5--6). Formulario: \fml{6-STG}, \fml{6-STM}, \fml{6-OPT}.}

% ---------------------------------------------------------------------
\question{A11}{Tsiolkovsky's law and the $\Delta V$ of a multistage rocket}
% source: cap06.tex l.21-110 (thrust, c, Isp), l.112-162 (Tsiolkovsky), l.164-354 (staging)
\begin{qbox}
State Tsiolkovsky's law and use it to evaluate the total $\Delta V$ of a multiple-stage rocket in terms of its propellant, inert,
and payload mass. Define the structural coefficient and specific impulse and discuss their effect on the total $\Delta V$.
\qmeta{Sep 2025 (listed as a 2025 question).}
\qnote{Adjacent to \qref{A10} (parameters of each stage) and \qref{A8} (losses with respect to this ideal law).}
\end{qbox}

\begin{outline}
\begin{enumerate}
\item Thrust and effective exhaust velocity $c=g_0I_{sp}$; $\dot m=-T/c$.
\item Assumptions; derivation $\dot v=-c\dot m/m$ $\Rightarrow$ $\Delta v=c\ln(m_0/m_f)$, $m_f=m_0e^{-\Delta v/c}$; meaning (upper bound).
\item Multistage: masses of each stage, $m_0^{(i)}$ and $m_f^{(i)}$ in terms of $m_p$, $m_s$, $m_{PL}$; $\Delta v_{tot}=\sum c_i\ln(m_0^{(i)}/m_f^{(i)})$.
\item Structural coefficient $\epsilon_s$ and specific impulse $I_{sp}$: definitions and effect on $\Delta v_{tot}$.
\end{enumerate}
\end{outline}

\answer
\paragraph{Thrust and specific impulse.} From the momentum balance of the vehicle and of the ejected mass,
$\underline T=\dot m\,\underline c_{ej}+A_e(p_e-p_a)\hat v=-\dot m\,(-\underline c)$, i.e.\ $T=|\dot m|\,c$, with the
\emph{effective exhaust velocity} $\underline c=\underline c_{ej}+A_e(p_e-p_a)\hat v/\dot m$ (includes the pressure thrust), and
\begin{equation}
\dot m=-\frac{T}{c},\qquad c=g_0I_{sp}.
\end{equation}
The \textbf{specific impulse} $I_{sp}$ [s] is the effective exhaust velocity divided by $g_0$ (sea-level gravity): the impulse
delivered per unit weight of propellant consumed. Typical: solid $\sim$270--290\,s, LOX/LH$_2$ $\sim$450\,s \notinnotes.

\paragraph{Tsiolkovsky's law.} Assumptions: no aerodynamic force, no gravity, thrust aligned with the velocity (valid for
impulsive manoeuvres or motion without $\underline G$ and $\underline A$). Then
\begin{equation}
\sys{\frac{dv}{dt}&=\frac{T}{m}\\ \dot m&=-\frac{T}{c}}\;\Rightarrow\;\dot v=-c\,\frac{\dot m}{m}
\;\Rightarrow\;\Delta v=v_f-v_0=c\ln\frac{m_0}{m_f},\qquad m_f=m_0\,e^{-\Delta v/c}.
\end{equation}
It relates the propellant consumption $m_p=m_0-m_f$ to the velocity change; it gives the \emph{maximum} $\Delta v$ for a given
mass ratio (minimum propellant for a given $\Delta v$): real ascents suffer gravity, drag and misalignment losses.

\paragraph{Multistage rocket.} Stage $j$ ($j=1\dots N$, burnt in this order) has propellant $m_p^{(j)}$ and inert mass $m_s^{(j)}$; the
payload is $m_{PL}$. Stage $i$ burns while all upper stages and the payload are attached and the lower stages have been jettisoned:
\begin{equation}
m_0^{(i)}=m_{PL}+\sum_{j=i}^{N}\left(m_s^{(j)}+m_p^{(j)}\right),\qquad m_f^{(i)}=m_0^{(i)}-m_p^{(i)},
\end{equation}
\begin{equation}
\boxed{\;\Delta v_{tot}=\sum_{i=1}^{N}c_i\ln\frac{m_0^{(i)}}{m_f^{(i)}}
=\sum_{i=1}^{N}g_0I_{sp}^{(i)}\ln\frac{m_{PL}+\sum_{j=i}^{N}\left(m_s^{(j)}+m_p^{(j)}\right)}{m_{PL}+m_s^{(i)}+\sum_{j=i+1}^{N}\left(m_s^{(j)}+m_p^{(j)}\right)}\;}
\end{equation}
In the notation of the notes, $m_u^{(i)}=m_0^{(i+1)}$, $u_0^{(i)}=m_f^{(i)}/m_0^{(i)}$ and $\Delta v_{tot}=-\sum_ic_i\ln u_0^{(i)}$.
The dropped inert mass $m_s^{(i)}$ is no longer accelerated by the upper stages: this is the gain of staging.

\paragraph{Structural coefficient and its effect.} $\epsilon_s^{(i)}=\dfrac{m_s^{(i)}}{m_s^{(i)}+m_p^{(i)}}$ (fraction of inert mass in
the stage; set by technology, $\epsilon_s\ge\epsilon_{s,\min}$). Since $u_0^{(i)}=\epsilon_s^{(i)}+(1-\epsilon_s^{(i)})\,m_u^{(i)}/m_0^{(i)}$, at
given payload fractions a smaller $\epsilon_s$ gives a smaller $u_0$ and a larger $\Delta v$; at given $\Delta v$, it gives a larger payload.
Each stage is limited by $\Delta v^{(i)}<c_i\ln(1/\epsilon_s^{(i)})$, and for similar optimal stages
$\Delta v_{\max}=-cN\ln\{\epsilon_s+(1-\epsilon_s)[m_u^{(N)}/m_0^{(1)}]^{1/N}\}$ decreases as $\epsilon_s$ increases.

\paragraph{Specific impulse and its effect.} $\Delta v$ is \emph{proportional} to $c=g_0I_{sp}$: at given masses, $\Delta v$ scales linearly
with $I_{sp}$ ($\partial\Delta v^{(i)}/\partial c_i=\ln(m_0^{(i)}/m_f^{(i)})$); at given $\Delta v$ the propellant fraction
$1-e^{-\Delta v/c}$ decreases exponentially with $I_{sp}$. High-$I_{sp}$ engines are therefore most valuable on the upper stages,
whose mass ratio multiplies the whole payload \notinnotes.

\begin{keyf}
$c=g_0I_{sp}$, $\dot m=-T/c$; \ $\Delta v=c\ln\dfrac{m_0}{m_f}$; \ $\Delta v_{tot}=\displaystyle\sum_i c_i\ln\frac{m_0^{(i)}}{m_0^{(i)}-m_p^{(i)}}$,
$m_0^{(i)}=m_{PL}+\displaystyle\sum_{j\ge i}(m_s^{(j)}+m_p^{(j)})$; \ $\epsilon_s=\dfrac{m_s}{m_s+m_p}$.
\end{keyf}

\pitfalls
\begin{itemize}
\item In stage $i$ the initial mass includes all upper stages and the payload, but not the jettisoned lower stages.
\item State the three assumptions of Tsiolkovsky's law and that it gives an upper bound.
\item $I_{sp}$ in seconds: multiply by $g_0=9.80665$\,m/s$^2$.
\end{itemize}

\source{cap06.tex l.21--110 (thrust, $c$, $I_{sp}$), l.112--162 (Tsiolkovsky), l.164--224 (staging parameters), l.225--354 (optimal staging).
Formulario: \fml{6-ISP}, \fml{6-TSI}, \fml{6-DVN}.}
